\subsubsection*{Ítem a}
A continuación se mostrará una matriz de confusión para la clasificación mediante LDA.
\begin{center}
    \begin{tabular}{|c|l|l|l|l|}
    \hline
    \multicolumn{2}{|l|}{\multirow{2}{*}{}}       & \multicolumn{3}{c|}{Clase predicha}            \\ \cline{3-5} 
    \multicolumn{2}{|l|}{}                        & Iris-setosa & Iris-versicolor & Iris-virginica \\ \hline
    \multirow{3}{*}{Clase real} & Iris-setosa     & 50          & 0               & 0              \\ \cline{2-5} 
                                & Iris-versicolor & 0           & 48              & 2              \\ \cline{2-5} 
                                & Iris-virginica  & 0           & 1               & 49             \\ \hline
    \end{tabular}
\end{center}

Tal como se observa en la presente tabla, el 98\% de los datos fueron predichos correctamente. Es decir, el error aparente es de $0.02$. 

Con validación cruzada (utilizamos leave one out) dio el mismo resultado.\\

\subsubsection*{Ítem b}
La matriz de confusión es idéntica a la obtenida en el ítem a, y por ende también el error aparente. Sin embargo, al realizar validación cruzada obtuvimos el siguiente resultado:

\begin{center}
    \begin{tabular}{|c|l|l|l|l|}
    \hline
    \multicolumn{2}{|l|}{\multirow{2}{*}{}}       & \multicolumn{3}{c|}{Clase predicha}            \\ \cline{3-5} 
    \multicolumn{2}{|l|}{}                        & Iris-setosa & Iris-versicolor & Iris-virginica \\ \hline
    \multirow{3}{*}{Clase real} & Iris-setosa     & 50          & 0               & 0              \\ \cline{2-5} 
                                & Iris-versicolor & 0           & 47              & 3             \\ \cline{2-5} 
                                & Iris-virginica  & 0           & 1               & 49             \\ \hline
    \end{tabular}
\end{center}

Siendo entonces el error por validación cruzada de $0.0267$.


